\begin{abstract}
Cross-domain few-shot object detection (CD-FSOD) adapts a pretrained detector to unseen domains using only a few labeled examples. Existing methods transfer visual prototypes, generated data, or multimodal prompts, but often treat every support factor as valid adaptation evidence. We identify \emph{support-induced context leakage}: changing only the support background can alter query classification and localization, even when the query image and support foreground remain fixed. We propose a knowledge-enhanced Hard-Soft framework that regulates both the evidence being transferred and its authority over model decisions. Foreground-Frozen Context Probing (FFCP) estimates class-wise context-sensitive subspaces through label-preserving support interventions. Context-Constrained Hard-Soft Decomposition (CHSD) constructs a Hard identity anchor and a reliability-weighted Soft target-domain prototype after removing context-sensitive directions. Asymmetric Task-Authority Routing (ATAR) uses Hard--Soft agreement as a reliability signal and applies a bounded cosine correction only to uncertain detector queries. The correction is zero at initialization and leaves the localization branch unchanged. The full model, denoted S2H-CD-FSOD, is evaluated on six CD-FSOD benchmarks under 1-, 5-, and 10-shot protocols. A controlled study on Clipart1k and NEU-DET separates context-induced category and box drift from training-seed variation, and a compact set of ablations and qualitative comparisons clarifies when support evidence helps and when its authority should be suppressed.
\end{abstract}
